\documentstyle[%


righttag%


,11pt%


,amssymb%


,russian%               remove in Eng.


,izvru%                 Set izven% in Eng.


]{amsart}





%


% verify \text before (13)


%\newcommand{\ov}{\overline}


%\newcommand{\wt}{\widehat}


\newcommand{\la}{\langle}


\newcommand{\ra}{\rangle}


\newcommand{\re}{\operatorname{Re}}


\newcommand{\im}{\operatorname{Im}}


\newcommand{\spaceL}{L^{\ov\alpha}_{p,r}}


\newcommand{\tts}[1]{}


\renewcommand{\baselinestretch}{1.05}


\renewcommand{\thesection}{\Roman{section}}





%\hoffset=-15mm%                  For Eng.





\setcounter{page}{24}


\god{1998}


\nomer{5~(432)} %                        Remove in Eng.


%\nomer{Vol.\,42, No.\,5}%               Set in Eng.


%\str{\pageref{begin}--\pageref{end}}%  Set in Eng.


\udk{517.993}


\author{\it А.Н.\,Карапетянц, В.А.\,Ногин}


% NB:   ^^ remove in Eng.


\title{Характеризация функций из анизотропных\\


пространств комплексного порядка}





\thanks{Работа поддержана Российским фондом


фундаментальных исследований (проект \No\,94-01-00577-A).}


\date{}


% \pagestyle{myheadings}%         Set in Eng.





\pagestyle{plain} %              Remove in Eng.


\begin{document}





% \thispagestyle{plain}%          Set in Eng.





\maketitle


\label{begin}





В работе получено описание анизотропных пространств


\begin{gather}


\spaceL\equiv \spaceL(\Bbb R^n)=\{f(x):\|f\|_{\spaceL}\equiv


\|f\|_r+\|F^{-1}(|\xi_1|^{\alpha_1}+\dotsb+|\xi_n|^{\alpha_n})


Ff\|_p<\infty\},\\


1<p<\infty,\ \;1\le r<\infty,\ \;


\ov\alpha=(\alpha_1,\dotsc,\alpha_n)\in\Bbb C^n,\ \; \re\alpha_j>0,


\ \;j=1,\dotsc,n,\notag


\end{gather}


$F$ --- преобразование Фурье, понимаемое в соответствующем смысле. При


вещественных $\alpha_j$ и $r=p$ пространства (1) совпадают с


введенными и


исследованными П.И.\,Лизоркиным (см. \cite{1}--\cite{4} и \cite{5})


пространствами $p$-суммируемых функций, имеющих в $L_p$ обобщенные


лиувиллевские производные $D^{\alpha_j}f$, $\alpha_j>0$, $j=1,\dotsc,n$,


по переменным $x_j$ соответственно.


В случае $r\ne p$ и $\alpha_1=\dotsb=\alpha_n=\alpha>0$


пространства $\spaceL$ были


изучены С.Г.\,Самко (см. \cite{6},


\cite{7} и \cite{8}, \cite{9}). В анизотропном случае пространства


$\spaceL$, а также пространства, совпадающие с ними с точностью до


эквивалентности норм, при вещественных $\alpha_j$ изучались в \cite{10},


\cite{11} (см. также \cite{12}--\cite{14}, где рассматривалась


анизотропность более общего характера). В указанных работах описание


класса $\spaceL$ было дано в терминах гиперсингулярных интегралов (ГСИ),


введенных П.И.\,Лизоркиным в \cite{3} (и их модификаций).





Интерес, проявленный нами к пространствам $\spaceL$,


$\alpha_j\in\Bbb C$, обусловлен, в


частности, тем, что они содержат классические пространства бесселевых,


риссовых, параболических потенциалов комплексного порядка.


Принципиальное отличие случая комплексных $\alpha_j$ состоит в том, что


для исследования пространств $\spaceL$ неприменимы, вообще говоря,


использовавшиеся в \cite{1}--\cite{11} и традиционные в таких вопросах


методы, базирующиеся на технике $p$-мультипликаторов. Это связано с


необходимостью деления на функцию


$|\xi_1|^{\alpha_1}+\dotsb+|\xi_n|^{\alpha_n}$,


которая в случае комплексных $\alpha_j$


таких, что


$\Delta_{jk}=\im\alpha_j\re\alpha_k-\im\alpha_k\re\alpha_j\ne 0$


хотя бы для двух индексов $j$ и $k$, имеет нули в $\Bbb R^n$,


отличные от начала координат (напр., при $n=2$ координаты нулей функции


$|\xi_1|^{\alpha_1}+|\xi_2|^{\alpha_2}$


в случае $\Delta_{12}\ne 0$ находятся из равенств:


$|\xi_1|=\exp\{\pi(2k+1)\re\alpha_2/\Delta_{12}\}$,


$|\xi_2|=\exp\{\pi(2k+1)\re\alpha_1/\Delta_{12}\}$, $k\in\Bbb Z$).


В частности, для описания классов $\spaceL$ неприменим метод ГСИ


(вообще не ясно, как


определить эти конструкции). Здесь используется другой


(аппроксимативный) подход, применявшийся ранее при обращении операторов


типа потенциала с $L_p$-плотностями (см. обзорную статью \cite{15} и


имеющуюся там библиографию). В рамках этого подхода описание дано в


терминах условно сходящихся (по $L_p$-норме) сверточных операторов с


суммируемыми ядрами, выражаемыми через элементарные функции (в отличие


от описания в терминах анизотропных ГСИ при вещественных $\alpha_j$; см.


ниже замечание 2).





Будем использовать обозначения:


$P_\varepsilon(x)=\frac{1}{\pi}


\frac{\varepsilon}{x^2+\varepsilon^2}$, $x\in\Bbb R^1$, --- ядро


Пуассона в $\Bbb R^1$;


$\la f,\omega\ra=\int\limits_{\Bbb R^n}\ov{f(x)}\omega(x)dx$;


$(Ff)(\xi)=\int\limits_{\Bbb R^n}f(x)e^{ix\xi}dx$


--- преобразование Фурье функции $f(x)$;


$(F^{-1}f)(x)=\frac{1}{(2\pi)^n}\int\limits_{\Bbb R^n}


f(\xi)e^{-ix\xi}d\xi$ --- обратное


преобразование


Фурье; $S=S(\Bbb R^n)$ --- пространство Л.\,Шварца быстро убывающих


гладких функций;


$\Psi_v$, $v\subset\Bbb R^n$, --- введенный и исследованный в


\cite{16}--\cite{19} (см. также


\cite{8}, \S\,3) класс шварцевых функций, исчезающих вместе со всеми своими


производными на множестве $v$; топология в пространстве $\Psi_v$


задается счетным набором согласованных сравнимых норм:


$$


\|\omega\|_N=


\sup\begin{Sb}|k|\le N \\ \xi\in\Bbb R^n\setminus v\end{Sb}


[\max\{(1+|\xi|^2)^{1/2},\,1/\rho(\xi,v)\}]^n


|(\cal D^k \omega)(\xi)|,\quad N=0,1,2,\dotsc,


$$


где $\rho(\xi,v)=\min\limits_{\eta\in v}|\xi-\eta|$;


$\Phi_v=F^{-1}(\Psi_v)$


--- пространство прообразов Фурье функций из $\Psi_v$. Если $v$ ---


совокупность координатных гиперплоскостей, то будем писать


$\Psi_v=\Psi$, $\Phi_v=\Phi$;


пространства $\Phi$, $\Psi$ были введены и изучены П.И.\,Лизоркиным


(\cite{1}--\cite{4}). Через $BC^m=BC^m(\Bbb R^n)$ будем обозначать класс


функций, ограниченных вместе со всеми своими производными до $m$-го порядка


включительно. Под $z^\alpha$, $\alpha\in\Bbb C$, будем понимать главную


ветвь ($k=0$) многозначной


функции $z^\alpha=\exp\{\alpha(\ln|z|+i\arg z+2\pi ki)\}$, являющуюся


аналитической


функцией в плоскости с разрезом по отрицательной вещественной полуоси.





\section{Описание пространств $\spaceL$ в терминах аппроксимативных


операторов}





Пространство $\spaceL$ определим соотношением (1), понимая в нем


преобразование


Фурье в смысле $\Psi'$-распределений. Такое определение корректно,


поскольку функция


$|\xi_1|^{\alpha_1}+\dotsb+|\xi_n|^{\alpha_n}$


является мультипликатором в $\Psi$. Для $\alpha_j\in\Bbb C$,


$\re\alpha_j>0$ и $\varepsilon>0$


положим


$$


R_{\ov\alpha,\varepsilon}(x)=(F^{-1}(


|\xi_1|^{\alpha_1}+\dotsb+|\xi_n|^{\alpha_n})


e^{-\varepsilon|\xi_1|-\dotsb-\varepsilon|\xi_n|})(x).


$$


В силу формулы


$$


R_{\alpha,\varepsilon}(|x|)\equiv\frac{1}{2\pi}


\int_{\Bbb R^1}|\xi|^\alpha e^{-\varepsilon|\xi|}


e^{-ix\xi}d\xi=(\Gamma(\alpha+1)/2\pi)


[(\varepsilon-i|x|)^{-\alpha-1}+(\varepsilon+i|x|


)^{-\alpha-1}]


$$


получаем


\begin{equation}


R_{\ov\alpha,\varepsilon}(x)=\sum^n_{j=1}


R_{\alpha_j,\varepsilon}(|x_j|)


\prod^n\begin{Sb}k=1\\ k\ne j\end{Sb}P_\varepsilon(x_k).


\end{equation}


Нетрудно показать, что


$R_{\ov\alpha,\varepsilon}(x)\in L_1(\Bbb R^n)$. Положим далее


\begin{equation}


(T^{\ov\alpha}_\varepsilon f)(x)=


(R_{\ov\alpha,\varepsilon}\ast f)(x).


\end{equation}





Основной результат статьи составляет





\begin{thm}


Пусть $1\le p\le 2$, $1\le r\le 2$, $\re\alpha_j>0$. Тогда


\begin{equation}


\spaceL=\{f(x):f\in L_r,\,


T^{\ov\alpha}f\in L_p\},


\end{equation}


где


\begin{equation}


T^{\ov\alpha}f=\overset{(L_p)}{\lim_{\varepsilon\to 0}}{}


(T^{\ov\alpha}_\varepsilon f)(x).


\end{equation}


Кроме того, $\|f\|_{\spaceL}=\|f\|_r+\|T^{\ov\alpha}f\|_p$.


\end{thm}





\begin{pf}


Пусть $f(x)\in\spaceL$. Покажем, что


\begin{equation}


(T^{\ov\alpha}_\varepsilon f)(x)=(P_\varepsilon\varphi)(x),\quad


\varphi=F^{-1}(|\xi_1|^{\alpha_1}+\dotsb+|\xi_n|^{\alpha_n})


Ff\in L_p,


\end{equation}


где


$(P_\varepsilon\varphi)(x)=\Bigl(\prod\limits^n_{k=1}


P_\varepsilon(t_k)\ast\varphi\Bigr)(x)$. Положим


$v=\{\xi\in\Bbb R^n:|\xi_1|^{\alpha_1}+\dotsb+|\xi_n|^{\alpha_n}\}$


и для $\omega\in\Phi_v$ рассмотрим функционал


$\la T^{\ov\alpha}_\varepsilon f,\omega\ra


=\la f,\ov{R_{\ov\alpha,\varepsilon}}\ast\omega\ra$. Заметим, что


$\ov{R_{\ov\alpha,\varepsilon}}\ast\omega\in\Phi_v$; это


следует из того,


что функция


$\ov{R_{\ov\alpha,\varepsilon}}(\xi)=


(|\xi_1|^{\ov\alpha_1}+\dotsb+|\xi_n|^{\ov\alpha_n})


e^{-\xi|\xi_1|-\dotsb-\xi|\xi_n|}$


является мультипликатором в $\Psi_v$. С учетом этого получаем


\begin{gather*}


\la T^{\ov\alpha}_\varepsilon f,\omega\ra=\frac{1}{(2\pi)^n}


\la\wt f,(|\xi_1|^{\ov\alpha_1}+\dotsb+|\xi_n|^{\ov\alpha_n})


e^{-\varepsilon|\xi_1|-\dotsb-\varepsilon|\xi_n|}


\wt\omega(\varepsilon)\ra= \\


=\frac{1}{(2\pi)^n}\la\wt\varphi,


e^{-\varepsilon|\xi_1|-\dotsb-\varepsilon|\xi_n|}


\wt\omega(\varepsilon)\ra=


\la\varphi,P_\varepsilon\omega\ra =


\la P_\varepsilon\varphi,\omega\ra.


\end{gather*}


Второе из равенств этой цепочки вытекает из соотношения


$\la\wt f,\wt\omega\ra=\la\wt\varphi,\wt\omega(\xi)/


(|\xi_1|^{\alpha_1}+\dotsb+|\xi_n|^{\alpha_n})\ra$, $\omega\in\Phi_v$.


Таким образом, пришли к равенству


\begin{equation}


\la T^{\ov\alpha}_\varepsilon f,\omega\ra =\la


P_\varepsilon\varphi,\omega\ra,


\end{equation}


справедливому для любой функции $\omega\in\Phi_v$. Покажем, что (7)


справедливо


для $\omega\in S$, откуда и будет следовать (6). Выберем


последовательность $\omega_k\in\Phi_v$,


аппроксимирующую функцию $\omega\in S$ по норме $L_{r'}$ и $L_{p'}$


одновременно. Возможность


такой аппроксимации функции из $S$ функциями из $\Phi_v$, где $v$ ---


произвольное множество нулевой меры и $r',p'\ge 2$, установлена в


\cite{17}--\cite{19} (в нашем случае соотношение $\text{mes\,}v=0$


очевидно). Переходя


в равенстве


$\la T^{\ov\alpha}_\varepsilon f,\omega_k\ra=


\la P_\varepsilon\varphi,\omega_k\ra$ к пределу при $k\to\infty$,


убеждаемся в


справедливости (7) на функциях из $S$. На основании (6) заключаем, что


$$


\overset{(L_p)}{\lim_{\varepsilon\to 0}}{}


T^{\ov\alpha}_\varepsilon f=\varphi\in L_p.


$$





Пусть теперь $f(x)\in L_r$,


$T^{\ov\alpha}f\in L_p$. Для $\omega\in\Phi$ имеем


\begin{multline}


\la T^{\ov\alpha}f,\omega\ra=\bigl\la


\overset{(L_p)}{\lim_{\varepsilon\to 0}}{}


T^{\ov\alpha}_\varepsilon f,\omega\bigr\ra=


\lim_{\varepsilon\to 0}\la T^{\ov\alpha}_\varepsilon f,\omega\ra=


\lim_{\varepsilon\to 0}\la f,\ov{R_{\ov\alpha,\varepsilon}}\ast


\omega\ra =\\


=\frac{1}{(2\pi)^n}\lim_{\varepsilon\to 0}


\la\wt f,(|\xi_1|^{\ov\alpha_1}+\dotsb+|\xi_n|^{\ov\alpha_n})


e^{-\varepsilon|\xi_1|-\dotsb-\varepsilon|\xi_n|}


\wt\omega(\xi)\ra= \\


=\frac{1}{(2\pi)^n}\la\wt f,


(|\xi_1|^{\ov\alpha_1}+\dotsb+|\xi_n|^{\ov\alpha_n})


\wt\omega(\xi)\ra=\la F^{-1}


(|\xi_1|^{\alpha_1}+\dotsb+|\xi_n|^{\alpha_n})Ff,\omega\ra.


\end{multline}


Второе из равенств этой цепочки следует из того, что сходимость в $L_p$


влечет сходимость в $\Phi'$, предпоследнее --- из того, что, как легко


показать,


$e^{-\varepsilon|\xi_1|-\dotsb-\varepsilon|\xi_n|}


\wt\omega(\xi)@>>{\varepsilon\to 0}>\wt\omega(\xi)$


в топологии $\Psi$. Таким образом,


$\la F^{-1}(|\xi_1|^{\alpha_1}+\dotsb+|\xi_n|^{\alpha_n})


Ff,\omega\ra=\la T^{\ov\alpha}f,\omega\ra$,


$\omega\in\Phi$, откуда следует, что


$f\in\spaceL$ и


$F^{-1}(|\xi_1|^{\alpha_1}+\dotsb+|\xi_n|^{\alpha_n})Ff=


T^{\ov\alpha}f$.


\end{pf}





Содержащееся в теореме 1 условие сходимости аппроксимативных операторов


можно заменить условием равномерной ограниченности $L_p$-норм интегралов


$T^{\ov\alpha}_\varepsilon f$. Именно, справедлива





\begin{thm}


Пусть $1<p\le 2$, $1\le r\le 2$, $\re\alpha_j>0$. Тогда


$\spaceL=\{f(x):f\in L_r,\,\sup\limits_{\varepsilon>0}


\|T^{\ov\alpha}_\varepsilon f\|_p<\infty\}$. При


этом $\|f\|_{\spaceL}=\|f\|_r+\sup\limits_{\varepsilon>0}


\|T^{\ov\alpha}_\varepsilon f\|_p$.


\end{thm}





\begin{pf}


Если $f\in\spaceL$, то, как было ранее показано, имеет место равенство


(6), в котором $\varphi=T^{\ov\alpha}f$. Отсюда следует


\begin{equation}


\sup_{\varepsilon>0}\|T^{\ov\alpha}_\varepsilon f\|_p\le


\|T^{\ov\alpha}f\|_p.


\end{equation}





Пусть $f\in L_r$ и


$\sup\limits_{\varepsilon>0}\|T^{\ov\alpha}_\varepsilon f\|_p<\infty$.


Тогда существует такая подпоследовательность $\varepsilon_k$, что


$W\text{-}\overset{(L_p)}{\lim\limits_{\varepsilon_k\to 0}}{}


T^{\ov\alpha}_\varepsilon f=\varphi\in L_p$,


где $W\text{-}\overset{(L_p)}{\lim}{}$


означает слабый предел в $L_p$. Заменив в цепочке выкладок (8)


$\overset{(L_p)}{\lim}{}$


на $W\text{-}\overset{(L_p)}{\lim}{}$,


получаем соотношение


$\la F^{-1}(|\xi_1|^{\alpha_1}+\dotsb+|\xi_n|^{\alpha_n})


Ff,\omega\ra=\la\varphi,\omega\ra$,


из которого следует, что $f\in\spaceL$.





Далее имеем


$\|T^{\ov\alpha}_\varepsilon f\|_p\le\sup\limits_{\varepsilon>0}


\|T^{\ov\alpha}_\varepsilon f\|_p$, откуда


предельным


переходом при $\varepsilon\to 0$ получаем


$\|T^{\ov\alpha}f\|_p\le\sup\limits_{\varepsilon>0}


\|T^{\ov\alpha}_\varepsilon f\|_p$. Отсюда и из (9)


следует 


$\|T^{\ov\alpha}f\|_p=\sup\limits_{\varepsilon>0}


\|T^{\ov\alpha}_\varepsilon f\|_p$.


\end{pf}





\begin{note}


В случае, когда $\alpha_j>0$, $j=1,\dotsc,n$, приведенные выше теоремы


справедливы при $1\le p,r<\infty$ ($p\ne 1$ в теореме 2).


Доказательство этого факта аналогично доказательству теорем 1, 2 с


учетом того, что пространство $\Phi_v=\Phi$, построенное по


совокупности координатных гиперплоскостей, плотно в $L_p$, $1<p<\infty$


(см. \cite{2}).


\end{note}





\begin{note}


Упомянутое выше описание пространств $\spaceL$ в терминах ГСИ имеет вид


\begin{equation}


\spaceL=\{f(x):f\in L_r,\,\Bbb D^{\ov\alpha}f\in L_p\},\quad


1<p<\infty,\ \; 1\le r<\infty,\ \;\alpha_j>0,


\end{equation}


где


\begin{equation}


\Bbb D^{\ov\alpha}f=\int_{\Bbb R^n}((\Delta^{2l}_t f)(x)/


\rho^{n+\alpha^*}(t))dt \equiv


\overset{(L_p)}{\lim_{\varepsilon\to 0}}{}


\int_{\rho(t)>\varepsilon}((\Delta^{2l}_t f)(x)/


\rho^{n+\alpha^*}(t))dt


\end{equation}


--- анизотропный ГСИ П.И.\,Лизоркина. В (11)\;


$\frac{1}{\alpha^*}=\frac{1}{n}\sum\limits^n_{i=1}


\frac{1}{\alpha_i}$, $\rho(t)$ --- неявная


функция, являющаяся решением уравнения


$\sum\limits^n_{i=1}x^2_i\rho^{-2\lambda_i}=1$,


$\lambda_i=\frac{\alpha_i}{\alpha^*}$,


$l=\frac{1}{2}\max\limits_{j\in\ov{1,n}}\alpha_j$. Содержащееся в


теореме 1 описание пространств $\spaceL$ (при вещественных $\alpha_j$)


является более


эффективным в том смысле, что ядра операторов


$T^{\ov\alpha}_\varepsilon$ выражаются


явно и через элементарные функции. Однако оно не обладает такой


конструктивностью, как описание (10): если $f(x)$ --- ``достаточно


хорошая'' функция, то предел в (11) равен интегралу по $\Bbb R^n$; в


то же время вопрос о том, существует ли предел


$\lim\limits_{\varepsilon\to 0}(T^{\ov\alpha}_\varepsilon f)(x)$ и


чему он равен, требует


специального исследования. Такое исследование проведено ниже в п.\,II, в


котором по существу речь идет о конструктивности описания (4).


\end{note}





\section{Сходимость аппроксимативных операторов на


``хороших'' функциях}





Здесь будет доказано существование предела


$\lim\limits_{\varepsilon\to 0}(T^{\ov\alpha}_\varepsilon f)(x)$,


$x\in\Bbb R^n$, на


функциях $f(x)$


класса $BC^m$,


$m\ge\max\{[\re\alpha_1],\dotsc,[\re\alpha_n]\}$. Для таких функций


$f(x)$, следуя (\cite{8}, с.\,86;


\cite{9}, (26.67)), введем регуляризацию расходящегося интеграла


$\int\limits_{\Bbb R^1}\frac{f(x-e_k t)}{|t|^{1+\alpha_k}}dt$, $e_k$


--- координатный орт, равенством


\begin{multline}


\text{p.\,f.\,}\int_{\Bbb R^1}


\frac{f(x-e_k t)}{|t|^{1+\alpha_k}}dt=\int_{\Bbb R^1}


\frac{f(x-e_k t)-\chi(t)P^{[\re\alpha_k]}_{t,k}(x)}


{|t|^{1+\alpha_k}}dt+ \\


+\sideset{}{'}\sum^{[\re\alpha_k]}_{s=0}


\frac{\Omega_s}{s!\,(s-\alpha_k)}


\frac{\partial^s}{\partial x^s_k}f(x),\quad


\Omega_s=


\begin{cases}


2, & s\text{ четное,}\\


0, & s\text{ нечетное},


\end{cases}


\end{multline}


где


$P^{[\re\alpha_k]}_{t,k}(x)=\sum\limits^{[\re\alpha_k]}_{s=0}


\frac{(-t)^s}{s!}\frac{\partial^s}{\partial x^s_k}f(x)$, $\chi(t)$ ---


характеристическая функция интервала $(-1,\,1)$, штрих у суммы


означает пропуск слагаемого с номером $s=\re\alpha_k$ в случае, когда


$\alpha_k$ --- целое число.





\begin{thm}


Пусть $f(x)\in BC^m$, $m\ge\max\{[\re\alpha_1],\dotsc,[\re\alpha_n]\}$.


Тогда предел


$\lim\limits_{\varepsilon>0}(T^{\ov\alpha}_\varepsilon f)(x)$


существует


для $\forall x\in\Bbb R^n$ и имеет место равенство


\begin{gather*}


\displaybreak[3]


\lim_{\varepsilon\to 0}(T^{\ov\alpha}_\varepsilon f)(x)=


\sum^n_{j=1}(J^{\alpha_j}f)(x),\\


J^{\alpha_j}f(x)=


\begin{cases}


\displaystyle


\frac{\Gamma(\alpha_j+1)\cos \frac{\pi(\alpha_j+1)}{2}}{\pi}


\text{\rom{ p.\,f. }}\int_{\Bbb R^1}


\frac{f(x-e_j t)}{|t|^{1+\alpha_j}}dt, & \alpha_j\ne 2,4,\dotsc;\\


\displaystyle\biggl(i\frac{\partial}{\partial x_j}


\biggr)^{\alpha_j}f(x), & \alpha_j=2,4,\dotsc


\end{cases}


\end{gather*}


\end{thm}





\begin{pf}


Интеграл $(T^{\ov\alpha}_\varepsilon f)(x)$ запишем в виде


\begin{gather}


(T^{\ov\alpha}_\varepsilon f)(x)=


\sum^n_{j=1}\frac{1}{\pi^{n-1}\varepsilon^{1+\alpha_j}}


\int_{\Bbb R^1}R_{\alpha_j}\Bigl(\frac{|t_j|}{\varepsilon}\Bigr)


dt_j\int_{\Bbb R^{n-1}}


\prod^n\begin{Sb}k=1\\ k\ne j\end{Sb}


\frac{\varepsilon}{t^2_k+\varepsilon^2}f(x-t)d\ov t_j=


\sum^n_{j=1}(J^{\alpha_j}_\varepsilon f)(x),\\


\ov t_j=(t_1,\dotsc,t_{j-1},\,t_{j+1},\dotsc,t_n),\notag


\end{gather}


где $R_{\alpha_j}(|t|)=R_{\alpha_j,1}(|t|)$. Покажем, что


$\lim\limits_{\varepsilon\to 0}(J^{\alpha_j}_\varepsilon f)(x)=


(J^{\alpha_j}f)(x)$, $x\in\Bbb R^n$,


откуда и будет следовать


утверждение теоремы. Пусть вначале $\alpha_j\ne 2,4,\dotsc$ Интеграл


$(J^{\alpha_j}_\varepsilon f)(x)$


представим в виде


\begin{multline*}


(J^{\alpha_j}_\varepsilon f)(x)=\\


=\frac{1}{\pi^{n-1}\varepsilon^{1+\alpha_j}}\int_{\Bbb R^1}


R_{\alpha_j}\Bigl(\frac{|t_j|}{\varepsilon}\Bigr)dt_j


\int_{\Bbb R^{n-1}}


\prod^n\begin{Sb}k=1\\ k\ne j\end{Sb}


\frac{1}{t^2_k+1}[f(x_j-t_j,\,\ov x_j-\varepsilon\ov t_j)


-\chi(t_j) P^{[\re\alpha_j]}_{t_j,j}


(x_j,\,\ov x_j-\varepsilon\ov t_j)]d\ov t_j+ \\


+\frac{1}{\pi^{n-1}\varepsilon^{1+\alpha_j}}


\int_{|t_j|<1}R_{\alpha_j}\Bigl(


\frac{|t_j|}{\varepsilon}\Bigr) dt_j


\int_{\Bbb R^{n-1}}


\prod^n\begin{Sb}k=1\\ k\ne j\end{Sb}


\frac{1}{t^2_k+1}


P^{[\re\alpha_j]}_{t_j,j}(x_j,\,\ov x_j-\varepsilon\ov t_j)


d\ov t_j\equiv


(J^{\alpha_j}_{1,\varepsilon}f)(x)+


(J^{\alpha_j}_{2,\varepsilon}f)(x),


\end{multline*}


где обозначено $f(x)=f(x_j,\,\ov x_j)$. Рассмотрим интеграл


$(J^{\alpha_j}_{1,\varepsilon}f)(x)$. Нетрудно


видеть, что


\begin{equation}


\frac{1}{\varepsilon^{1+\alpha_j}}R_{\alpha_j}\Bigl(


\frac{|t_j|}{\varepsilon}\Bigr)@>>{\varepsilon\to 0}>


\frac{\Gamma(\alpha_j+1)\cos


\frac{\pi(\alpha_j+1)}{2}}{\pi}


|t_j|^{-1-\alpha_j}


\end{equation}


и, кроме того,


\begin{multline}


\biggl|\varepsilon^{-1-\alpha_j}R_{\alpha_j}\Bigl(


\frac{|t_j|}{\varepsilon}\Bigr)\int_{\Bbb R^{n-1}}


\prod^n\begin{Sb}k=1\\ k\ne j\end{Sb}


\frac{1}{t^2_k+1}[f(x_j-t_j,\,\ov x_j-\varepsilon\ov t_j)-


\chi(t_j)P^{[\re\alpha_j]}_{t_j,j}(x_j,\,\ov x_j-


\varepsilon\ov t_j)]dt_j\biggr|\le \\


\le


C


\begin{cases}


|t_j|^{-1-\re\alpha_j}, & |t_j|>1;\\


|t_j|^{-\{\re\alpha_j\}}, & |t_j|\le 1,


\end{cases}


\end{multline}


где $C$ не зависит от $\varepsilon$.


Переходя на основании (15) к пределу


под знаком интеграла (по $dt_j$) в


$(J^{\alpha_j}_{1,\varepsilon}f)(x)$, с учетом (14) получаем


\begin{equation}


\lim_{\varepsilon\to 0}(J^{\alpha_j}_{1,\varepsilon}f)(x)=


\frac{\Gamma(\alpha_j+1)\cos\frac{\pi(\alpha_j+1)}{2}}{\pi}


\int_{\Bbb R^1}


\frac{f(x-e_j t)-\chi(t_j)


P^{[\re\alpha_j]}_{t_j,j}(x)}


{|t_j|^{1+\alpha_j}}dt_j.


\end{equation}





Обратимся теперь к $(J^{\alpha_j}_{2,\varepsilon}f)(x)$. Если


$\re\alpha_j\ne 2,4,\dotsc,$ то в силу


$\int\limits_{\Bbb R^1}t^k_j R_{\alpha_j}(|t_j|)dt_j=0$ для


$k<\re\alpha_j$ и


$\int\limits_{|t_j|<1}t^{\re\alpha_j}_j R_{\alpha_j}\Bigl(


\frac{|t_j|}{\varepsilon}\Bigr)dt_j=0$


при нечетных $\re\alpha_j$ получаем


$$


(J^{\alpha_j}_{2,\varepsilon}f)(x)=


-\frac{1}{\pi^{n-1}}\sideset{}{'}\sum^{[\re\alpha_j]}_{s=0}


\int_{|t_j|>1}\frac{(-t)_j^s}{s!\,\varepsilon^{1+\alpha_j}}


R_{\alpha_j}\Bigl(\frac{|t_j|}{\varepsilon}\Bigr)dt_j


\int_{\Bbb R^{n-1}}\prod^n\begin{Sb}k=1\\ k\ne j\end{Sb}


\frac{1}{t^2_k+1}\frac{\partial^s}{\partial x^s_j}


f(x_j,\,\ov x_j-\varepsilon\ov t_j)d\ov t_j.


$$


Переходя в правой части последнего равенства к пределу под


знаком интеграла, с учетом (14) имеем


\begin{equation}


\lim_{\varepsilon\to 0}(J^{\alpha_j}_{2,\varepsilon}f)(x)=


-\frac{\Gamma(\alpha_j+1)\cos\frac{\pi(\alpha_j+1)}{2}}{\pi}\;


\sideset{}{'}\sum^{[\re\alpha_j]}_{s=0}


\frac{\Omega_s}{s!\,(\alpha_j-s)}


\frac{\partial^s}{\partial x^s_j}f(x).


\end{equation}





Если же $\re\alpha_j=2,4,\dotsc,$ $\im\alpha_j\ne 0$,


то, интегрируя по частям, имеем


$$


\lim_{\varepsilon\to 0}\frac{1}{(\re\alpha_j)!}


\int_{|t_j|<1}t^{\re\alpha_j}R_{\alpha_j,\varepsilon}


(|t_j|)dt_j=-


\frac{2\Gamma(\alpha_j+1)\cos\frac{\pi(\alpha_j+1)}{2}}


{\pi i\im\alpha_j(\re\alpha_j)!}


\biggl(\frac{\partial}{\partial x_j}\biggr)^{\re\alpha_j}f(x),


$$


откуда, как и при $\re\alpha_j\ne 2,4,\dotsc$, получаем (17). Из (16),


(17) следует утверждение теоремы при $\alpha_j\ne 2,4,\dotsc$





Остается рассмотреть случай $\alpha_j=2,4,\dotsc$ Проинтегрировав


$(J^{\alpha_j}_\varepsilon f)(x)$ из (13) по


частям $\alpha_j$ раз, будем иметь


$$


(J^{\alpha_j}_\varepsilon f)(x)=\frac{1}{\pi^n}


\int_{\Bbb R^n}\prod^n_{k=1}\frac{1}{t^2_k+1}


\biggl(\Bigl(i\frac{\partial}{\partial x_j}\Bigr)^{\alpha_j}f


\biggr)(x-\varepsilon t)dt.


$$


Переходя в этом равенстве к пределу при $\varepsilon\to 0$, получаем


$\lim\limits_{\varepsilon\to 0}(J^{\alpha_j}_\varepsilon f)(x)=


\Bigl(\Bigl(i\frac{\partial}{\partial x_j}\Bigr)^{\alpha_j}f


\Bigr)(x)$.


\end{pf}








\begin{thebibliography}{99}





\bibitem{1}


Лизоркин П.И. {\em Обобщенное лиувиллевское дифференцирование и


функциональные пространства $L^r_p(E_n)$. Теоремы вложения} //


Матем. сб. -- 1963. -- Т.\,60. -- \No\,3. -- С.\,325--353.


\bibitem{2}


Лизоркин П.И. {\em Обобщенное лиувиллевское дифференцирование и метод


мультипликаторов в теории вложений классов дифференцируемых функций} //


Тр. Матем. ин-та АН СССР. -- 1969. -- Т.\,105. -- С.\,89--167.


\bibitem{3}


Лизоркин П.И. {\em Описание пространств $L^r_p(R^n)$ в терминах


разностных сингулярных интегралов} // Матем. сб. -- 1970. -- Т.\,81. --


\No\,1. -- С.\,79--91.


\bibitem{4}


Лизоркин П.И. {\em Операторы, связанные с дробным дифференцированием, и


классы дифференцируемых функций} // Тр. Матем. ин-та АН СССР. -- 1972.


-- Т.\,117. -- С.\,212--243.


\bibitem{5}


Никольский С.М. {\em Приближение функций многих переменных и теоремы


вложения}. -- М.: Наука, 1983. -- 480~с.


\bibitem{6}


Самко С.Г. {\em О пространствах риссовых потенциалов} // Изв. АН СССР.


Сер. матем. -- 1976. -- Т.\,40.-- \No\,5. -- С.\,1143--1172.


\bibitem{7}


Самко С.Г. {\em Пространства $L^\alpha_{p,r}(\Bbb R^n)$ и


гиперсингулярные интегралы} // Studia math. -- 1977. -- T.\,61. --


\No\,3. -- S.\,193--220.


\bibitem{8}


Самко С.Г. {\em Гиперсингулярные интегралы и их приложения}. --


Ростов-на-Дону: Изд-во Рост. ун-та, 1984. -- 208~с.


\bibitem{9}


Самко С.Г., Килбас А.А., Маричев О.И. {\em Интегралы и производные


дробного порядка и некоторые их приложения}. -- Минск: Наука и техника,


1987. -- 688~с.


\bibitem{10}


Давтян А.А. {\em Пространства анизотропных потенциалов. Приложения} //


Тр. Матем. ин-та АН СССР. -- 1986. -- Т.\,173. -- С.\,113--124.


\bibitem{11}


Емгушева Г.П., Ногин В.А. {\em Характеризация функций из анизотропных


классов лиувиллевского типа} // Изв. вузов. Математика. -- 1989. --


\No\,7. -- С.\,63--66.


\bibitem{12}


Dappa H., Trebels W. {\em On $L^1$ criteria for quasi-radial Fourier


multipliers with applications to some anisotropic function spaces} //


Anal. math. -- 1983. -- V.\,9. -- \No\,4. -- P.\,275--289.


\bibitem{13}


Dappa H., Trebels W. {\em On hypersingular integrals and anisotropic


Bessel potential spaces} // Trans. Amer. Math. Soc. -- 1984. --


V.\,286. -- \No\,1. -- P.\,419--429.


\bibitem{14}


Хамидова Т.А. {\em Гиперсингулярные интегралы и пространства анизотропных


риссовых потенциалов}. -- Ростовск. ун-т. -- Ростов-на-Дону,


1988. -- 22 с. -- Деп. в ВИНИТИ 27.12.88, \No\,9067-B.


\bibitem{15}


Samko S.G. {\em Inversion theorems for potential type integral


transforms in $\Bbb R^n$ and on $S^{n-1}$} // Integral transforms and


special functions. -- 1993. -- V.\,1. -- \No\,2. -- P.\,145--163.


\bibitem{16}


Samko S.G. {\em Denseness of the spaces $\Phi_v$ of Lizorkin type in


the mixed $L_{\ov p}(\Bbb R^n)$-spaces} // Studia math. -- 1995.


-- T.\,113. -- \No\,3.


-- S.\,199--210.


\bibitem{17}


Самко С.Г. {\em Об основных функциях, исчезающих на заданном множестве,


и о делении на функции} // Матем. заметки. -- 1977. -- Т.\,21. --


\No\,5. -- С.\,677--689.


\bibitem{18}


Самко С.Г. {\em О плотности в $L_p(\Bbb R^n)$ пространств $\Phi_v$ типа


Лизоркина} // Матем. заметки. -- 1982. -- Т.\,31. -- \No\,6. --


С.\,855--865.


\bibitem{19}


Самко С.Г. {\em О плотности пространств $\Phi_v$ типа Лизоркина в


пространствах $L_{\ov p}(\Bbb R^n)$ со смешанной нормой} // Докл. РАН. --


1991. -- Т.\,319. -- \No\,3. -- С.\,567--569.








\end{thebibliography}





\vskip 2cm





\post{\it Ростовский государственный}{\it Поступила}


\post{\it университет}{10.06.1996}





\label{end}


\end{document}


